\documentclass[10pt,a4paper]{amsart}
\usepackage{amsmath,amssymb,amsthm,graphicx}
\usepackage{verbatim}
\newenvironment{reference}{\comment}{\endcomment}
%\newenvironment{reference}{}{}
\newenvironment{slogan}{\comment}{\endcomment}
\newenvironment{history}{\comment}{\endcomment}
% For commutative diagrams we use Xy-pic
\usepackage[all]{xy}

% We use 2cell for 2-commutative diagrams.
\xyoption{2cell}
% Theorem environments.
%
\theoremstyle{plain}
\newtheorem{theorem}[subsection]{Theorem}
\newtheorem{proposition}[subsection]{Proposition}
\newtheorem{lemma}[subsection]{Lemma}

\theoremstyle{definition}
\newtheorem{definition}[subsection]{Definition}
\newtheorem{example}[subsection]{Example}
\newtheorem{exercise}[subsection]{Exercise}
\newtheorem{situation}[subsection]{Situation}

\theoremstyle{remark}
\newtheorem{remark}[subsection]{Remark}
\newtheorem{remarks}[subsection]{Remarks}

\numberwithin{equation}{subsection}

% Macros
%
\def\lim{\mathop{\mathrm{lim}}\nolimits}
\def\colim{\mathop{\mathrm{colim}}\nolimits}
\def\Spec{\mathop{\mathrm{Spec}}}
\def\Hom{\mathop{\mathrm{Hom}}\nolimits}
\def\Ext{\mathop{\mathrm{Ext}}\nolimits}
\def\SheafHom{\mathop{\mathcal{H}\!\mathit{om}}\nolimits}
\def\SheafExt{\mathop{\mathcal{E}\!\mathit{xt}}\nolimits}
\def\Sch{\mathit{Sch}}
\def\Mor{\mathop{\mathrm{Mor}}\nolimits}
\def\Ob{\mathop{\mathrm{Ob}}\nolimits}
\def\Sh{\mathop{\mathit{Sh}}\nolimits}
\def\NL{\mathop{N\!L}\nolimits}
\def\CH{\mathop{\mathrm{CH}}\nolimits}
\def\proetale{{pro\text{-}\acute{e}tale}}
\def\etale{{\acute{e}tale}}
\def\QCoh{\mathit{QCoh}}
\def\Ker{\mathop{\mathrm{Ker}}}
\def\Im{\mathop{\mathrm{Im}}}
\def\Coker{\mathop{\mathrm{Coker}}}
\def\Coim{\mathop{\mathrm{Coim}}}

% Boxtimes
%
\DeclareMathSymbol{\boxtimes}{\mathbin}{AMSa}{"02}

%
% Macros for moduli stacks/spaces
%
\def\QCohstack{\mathcal{QC}\!\mathit{oh}}
\def\Cohstack{\mathcal{C}\!\mathit{oh}}
\def\Spacesstack{\mathcal{S}\!\mathit{paces}}
\def\Quotfunctor{\mathrm{Quot}}
\def\Hilbfunctor{\mathrm{Hilb}}
\def\Curvesstack{\mathcal{C}\!\mathit{urves}}
\def\Polarizedstack{\mathcal{P}\!\mathit{olarized}}
\def\Complexesstack{\mathcal{C}\!\mathit{omplexes}}
% \Pic is the operator that assigns to X its picard group, usage \Pic(X)
% \Picardstack_{X/B} denotes the Picard stack of X over B
% \Picardfunctor_{X/B} denotes the Picard functor of X over B
\def\Pic{\mathop{\mathrm{Pic}}\nolimits}
\def\Picardstack{\mathcal{P}\!\mathit{ic}}
\def\Picardfunctor{\mathrm{Pic}}
\def\Deformationcategory{\mathcal{D}\!\mathit{ef}}
\usepackage{cleveref}
\crefname{originallemma}{original lemma}{original lemmas}
\Crefname{originallemma}{Original lemma}{Original lemmas}
\makeatletter
\newlabel{lemma-adjoint}{{lemma-adjoint}{source}{}{sourcepointer-1}{}}
\newlabel{lemma-adjoint@cref}{{[originallemma][0][]lemma-adjoint}{[1][0][]source}{}{}{}}
\makeatother
\crefname{originallemma}{original lemma}{original lemmas}
\Crefname{originallemma}{Original lemma}{Original lemmas}
\makeatletter
\newlabel{lemma-local-cohomology-adjoint}{{lemma-local-cohomology-adjoint}{source}{}{sourcepointer-2}{}}
\newlabel{lemma-local-cohomology-adjoint@cref}{{[originallemma][0][]lemma-local-cohomology-adjoint}{[1][0][]source}{}{}{}}
\makeatother
\crefname{originalreference}{original reference}{original references}
\Crefname{originalreference}{Original reference}{Original references}
\makeatletter
\newlabel{more-algebra-lemma-divide-by-torsion}{{more-algebra-lemma-divide-by-torsion}{source}{}{sourcepointer-3}{}}
\newlabel{more-algebra-lemma-divide-by-torsion@cref}{{[originalreference][0][]more-algebra-lemma-divide-by-torsion}{[1][0][]source}{}{}{}}
\makeatother
\title{Dividing a module by its torsion in the ideal makes the local cohomology concentrate in degree zero}
\author{The Stacks Project Authors}
\date{stacks-project:dualizing.tex:local-snapshot-20261012; GNU FDL 1.2 or later}
\begin{document}
\maketitle
\noindent\textit{Editorial note.} All delivered mathematical text is the author's own, copied from the pinned Stacks Project chapter source sources/stacks-dualizing/pinned-original-source.tex; no statement or proof step was rewritten, reordered or removed.\par\medskip
\noindent\textit{Editorial note.} Preamble provenance: the Stacks Project keeps one single shared preamble for every chapter; this card takes it directly from the staged sources/stacks-dualizing/preamble.tex as three verbatim spans (lines 8-12, 14-18 and 154-217), so no derived file was written and nothing was added or redefined. The wrapper class and the base packages are supplied by the assembler.\par\medskip
\noindent\textit{Editorial note.} Bibliography provenance: neither the counted statement nor its proof contains a citation (this lane rescans the whole span before writing the card), so the card carries no bibliography and no entry was reconstructed.\par\medskip
\noindent\textit{Editorial note.} Selection provenance: this card was opened under the file-based control contract (CONTROL.md), with the mechanical pre-card gate run on its statement and proof interval before the card was written.\par\medskip
\noindent\textit{Allowed standard prerequisites (editorial).}\begin{itemize}
\item The identification of the zeroth local cohomology with support in a finitely generated ideal with the kernel of the map to the localizations. Invoked at the start of the counted proof (source lines 2127-2131). Computes the zeroth cohomology with support and proves the equality with the closed subscheme version.
\item The vanishing of the higher local cohomology of a torsion module. Invoked inside the counted proof (source lines 2135-2136). Gives the vanishing of the higher cohomology of the torsion submodule.
\item The characterization of modules without torsion as those annihilated by the functor from the proof of the first statement. Invoked inside the counted proof (source lines 2139-2140). Shows that the quotient by the torsion has no zeroth local cohomology.
\end{itemize}\medskip
\ifdefined\hypertarget\hypertarget{sourcepointer-1}{}\fi
\noindent\textit{Source reference.} Original lemma lemma-adjoint: sources/stacks-dualizing/pinned-original-source.tex, source lines 1294--1294. Navigation pointer only: the counted text refers to the Stacks item lemma-adjoint, which is an existing labeled result of this same chapter and is not part of this card. The pointer carries the target's own label as its number and its definition line as its span; its content is not reproduced here.\par\medskip
\ifdefined\hypertarget\hypertarget{sourcepointer-2}{}\fi
\noindent\textit{Source reference.} Original lemma lemma-local-cohomology-adjoint: sources/stacks-dualizing/pinned-original-source.tex, source lines 1466--1466. Navigation pointer only: the counted text refers to the Stacks item lemma-local-cohomology-adjoint, which is an existing labeled result of this same chapter and is not part of this card. The pointer carries the target's own label as its number and its definition line as its span; its content is not reproduced here.\par\medskip
\ifdefined\hypertarget\hypertarget{sourcepointer-3}{}\fi
\noindent\textit{Source reference.} Original reference more-algebra-lemma-divide-by-torsion: sources/stacks-dualizing/pinned-original-source.tex, source lines 2140--2140. Navigation pointer only: the counted text refers to the Stacks item more-algebra-lemma-divide-by-torsion, which belongs to another chapter of the project and is not part of this card. The pointer carries the target's own label as its number because this chapter's staged snapshot does not contain the other chapter's tag table; its content is not reproduced here.\par\medskip
\begin{lemma}
\label{lemma-divide-by-torsion}
Let $A$ be a ring and let $I \subset A$ be a finitely generated ideal.
Let $M$ be an $A$-module. Let $Z = V(I)$.
Then $H^0_I(M) = H^0_Z(M)$. Let $N$ be the common value and
set $M' = M/N$. Then
\begin{enumerate}
\item $H^0_I(M') = 0$ and $H^p_I(M) = H^p_I(M')$ and $H^p_I(N) = 0$
for all $p > 0$,
\item $H^0_Z(M') = 0$ and $H^p_Z(M) = H^p_Z(M')$ and $H^p_Z(N) = 0$
for all $p > 0$.
\end{enumerate}
\end{lemma}

\begin{proof}
By definition $H^0_I(M) = M[I^\infty]$ is $I$-power torsion.
By Lemma \ref{lemma-local-cohomology-adjoint} we see that
$$
H^0_Z(M) = \Ker(M \longrightarrow M_{f_1} \times \ldots \times M_{f_r})
$$
if $I = (f_1, \ldots, f_r)$. Thus $H^0_I(M) \subset H^0_Z(M)$ and
conversely, if $x \in H^0_Z(M)$, then it is annihilated by a $f_i^{e_i}$
for some $e_i \geq 1$ hence annihilated by some power of $I$.
This proves the first equality and moreover $N$ is $I$-power torsion.
By Lemma \ref{lemma-adjoint} we see that $R\Gamma_I(N) = N$.
By Lemma \ref{lemma-local-cohomology-adjoint} we see that $R\Gamma_Z(N) = N$.
This proves the higher vanishing of $H^p_I(N)$ and $H^p_Z(N)$ in (1) and (2).
The vanishing of $H^0_I(M')$ and $H^0_Z(M')$ follow from the preceding
remarks and the fact that $M'$ is $I$-power torsion free by
More on Algebra, Lemma \ref{more-algebra-lemma-divide-by-torsion}.
The equality of higher cohomologies for $M$ and $M'$ follow
immediately from the long exact cohomology sequence.
\end{proof}

\end{document}
